% ============================================================
% DOCUMENTCLASS
% ============================================================
\documentclass[12pt,oneside]{book}

% ============================================================
% METADATOS
% ============================================================
\newcommand{\universidad}{Universidad de Buenos Aires (UBA)}
\newcommand{\materia}{Derecho Procesal Civil y Comercial}
\newcommand{\titulo}{Prueba Testimonial y Prueba de Confesión}
\newcommand{\autor}{Isaac Edgar Camacho Ocampo}
\newcommand{\anio}{2024}
\newcommand{\reflexionportada}{Comprender los fundamentos procesales de la prueba es el primer paso para argumentar con solidez ante los tribunales.}

% ============================================================
% PAQUETES
% ============================================================
\usepackage[utf8]{inputenc}
\usepackage[T1]{fontenc}
\usepackage[spanish,es-nodecimaldot]{babel}

\usepackage[
    top=2.5cm,
    bottom=2.5cm,
    left=3cm,
    right=2.5cm,
    headheight=15pt
]{geometry}

\usepackage{amsmath}
\usepackage{amssymb}
\usepackage{mathtools}

\usepackage{graphicx}
\usepackage{float}
\usepackage{caption}
\usepackage{subcaption}

\usepackage{booktabs}
\usepackage{array}
\usepackage{longtable}
\usepackage{tabularx}

\usepackage{enumitem}

\usepackage{xcolor}
\usepackage[most]{tcolorbox}

\usepackage{fancyhdr}
\usepackage{ragged2e}

% ============================================================
% CONFIGURACIÓN
% ============================================================
\graphicspath{
    {./img/}
    {./graficos/}
    {./figuras/}
}

\setcounter{tocdepth}{2}

\setlength{\parindent}{0pt}
\setlength{\parskip}{0.5em}

\setlist[itemize]{
    topsep=0.3em,
    itemsep=0.2em
}

\setlist[enumerate]{
    topsep=0.3em,
    itemsep=0.2em
}

\clubpenalty=10000
\widowpenalty=10000

% ============================================================
% COMANDOS
% ============================================================
\newcommand{\abs}[1]{\left|#1\right|}
\newcommand{\norm}[1]{\left\lVert#1\right\rVert}
\newcommand{\vect}[1]{\mathbf{#1}}
\newcommand{\deriv}[2]{\frac{d#1}{d#2}}
\newcommand{\pderiv}[2]{\frac{\partial#1}{\partial#2}}

% ============================================================
% ENTORNOS / CAJAS
% ============================================================
\newtcolorbox{definition}[1]{
    enhanced,
    breakable,
    title={Definición: #1},
    fonttitle=\bfseries,
    colback=blue!3,
    colframe=blue!50!black,
    boxrule=0.8pt,
    arc=2mm
}

\newtcolorbox{important}{
    enhanced,
    breakable,
    title={Importante},
    fonttitle=\bfseries,
    colback=yellow!8,
    colframe=orange!70!black,
    boxrule=0.8pt,
    arc=2mm
}

\newtcolorbox{example}[1]{
    enhanced,
    breakable,
    title={Ejemplo: #1},
    fonttitle=\bfseries,
    colback=green!4,
    colframe=green!40!black,
    boxrule=0.8pt,
    arc=2mm
}

\newtcolorbox{formula}[1]{
    enhanced,
    breakable,
    title={Fórmula / Regla: #1},
    fonttitle=\bfseries,
    colback=purple!4,
    colframe=purple!50!black,
    boxrule=0.8pt,
    arc=2mm
}

\newtcolorbox{exercise}[1]{
    enhanced,
    breakable,
    title={Ejercicio: #1},
    fonttitle=\bfseries,
    colback=gray!5,
    colframe=gray!60!black,
    boxrule=0.8pt,
    arc=2mm
}

\newtcolorbox{note}{
    enhanced,
    breakable,
    title={Nota},
    fonttitle=\bfseries,
    colback=white,
    colframe=gray!50,
    boxrule=0.6pt,
    arc=2mm
}

% ============================================================
% CONFIGURACIÓN FINAL (encabezados, hipervínculos)
% ============================================================
\usepackage[
    colorlinks=true,
    linkcolor=blue!50!black,
    citecolor=green!40!black,
    urlcolor=blue!60!black,
    pdftitle={\titulo},
    pdfauthor={\autor},
    pdfsubject={Apuntes universitarios},
    bookmarks=true
]{hyperref}

\pagestyle{fancy}
\fancyhf{}

\fancyhead[L]{\nouppercase{\leftmark}}
\fancyhead[R]{\materia}

\fancyfoot[C]{\thepage}

\renewcommand{\headrulewidth}{0.4pt}
\renewcommand{\footrulewidth}{0pt}

\fancypagestyle{plain}{
    \fancyhf{}
    \fancyfoot[C]{\thepage}
    \renewcommand{\headrulewidth}{0pt}
}

% ============================================================
% DOCUMENT
% ============================================================
\begin{document}

% ------------------------------------------------------------
% PORTADA
% ------------------------------------------------------------
\begin{titlepage}

\thispagestyle{empty}

\begin{center}

\vspace*{1cm}

{\large \textsc{\universidad}}

\vspace{3cm}

{\Huge \textbf{\materia}}

\vspace{1cm}

{\LARGE \titulo}

\vspace{0.8cm}

\textit{\reflexionportada}

\vspace{1.2cm}

\rule{0.70\textwidth}{0.5pt}

\vspace{0.5cm}

{\large Apuntes de estudio}

\vspace{0.5cm}

\rule{0.70\textwidth}{0.5pt}

\vfill

{\large \autor}

\vspace{0.3cm}

{\large \anio}

\end{center}

\vfill

\begin{flushleft}
\textit{Este es un modesto aporte a los estudiantes de ``\materia'' no pretende sustituir el material oficial de la materia.}
\end{flushleft}

\end{titlepage}

% ------------------------------------------------------------
% ÍNDICE
% ------------------------------------------------------------
\tableofcontents

% ------------------------------------------------------------
% INTRODUCCIÓN
% ------------------------------------------------------------
\chapter*{Introducción}
\addcontentsline{toc}{chapter}{Introducción}

Esta unidad aborda dos medios de prueba fundamentales del proceso civil y comercial argentino: la prueba testimonial (declaración de testigos) y la prueba de confesión (absolución de posiciones). En ambos casos se analiza el régimen legal aplicable conforme al Código Procesal Civil y Comercial de la Nación (CPCCN), los deberes y derechos de los declarantes, las formas válidas de ofrecimiento, producción y valoración de la prueba, y los supuestos de caducidad que operan objetivamente ante la inactividad procesal de las partes. Se incluyen además referencias al Código Civil y Comercial de la Nación, jurisprudencia de la Corte Suprema de Justicia de la Nación y ejemplos prácticos del ejercicio forense.

Estos contenidos pueden comprenderse a partir de una analogía sencilla: el proceso judicial resulta comparable a una investigación periodística. Primero es necesario obtener los testimonios de quienes presenciaron los hechos (prueba testimonial), y luego lograr reconocimientos de los propios involucrados (prueba de confesión). Cada uno de estos medios de prueba, a su vez, atraviesa tres etapas comunes: ofrecimiento, producción y valoración.

\section*{Esquema temático de la unidad}
\addcontentsline{toc}{section}{Esquema temático de la unidad}

\begin{longtable}{
    >{\raggedright\arraybackslash}p{0.30\textwidth}
    >{\raggedright\arraybackslash}p{0.63\textwidth}
}
\caption{Prueba testimonial (arts. 426 a 456 CPCCN)} \\
\toprule
\textbf{Tema} & \textbf{Palabras clave} \\
\midrule
\endfirsthead
\toprule
\textbf{Tema} & \textbf{Palabras clave} \\
\midrule
\endhead
Concepto de testigo & tercero, declaración, hechos, sentidos \\
Objeto del testimonio & hechos alegados, valor según percepción \\
Deberes del testigo & comparecer, declarar, decir verdad \\
Testigos excluidos --- art. 427 & consanguíneos, afines, cónyuge, excepción familia \\
Tres momentos: ofrecimiento & datos, objeto, hasta ocho testigos, art. 333 \\
Producción: citación y audiencia & cédula, tres días, audiencia principal y supletoria \\
Interrogatorio preliminar --- art. 441 & generales de la ley, parentesco, interés \\
Forma de las preguntas --- art. 443 & hecho único, claras, no afirmativas, técnicas \\
Testigo técnico vs. perito & fungibilidad, conocimiento previo, designación \\
Negativa a responder --- art. 444 & secreto profesional, honor, autoincriminación \\
Caducidad --- arts. 432 y 437 & cuatro supuestos objetivos, desistimiento \\
Idoneidad del testigo --- art. 456 & incidente, sana crítica, impugnación \\
\bottomrule
\end{longtable}

\begin{longtable}{
    >{\raggedright\arraybackslash}p{0.30\textwidth}
    >{\raggedright\arraybackslash}p{0.63\textwidth}
}
\caption{Prueba de confesión (arts. 404 a 424 CPCCN)} \\
\toprule
\textbf{Tema} & \textbf{Palabras clave} \\
\midrule
\endfirsthead
\toprule
\textbf{Tema} & \textbf{Palabras clave} \\
\midrule
\endhead
Concepto de confesión & declaración de parte, efecto desfavorable, confesión provocada \\
Formas: espontánea y provocada & contestación de demanda, absolución de posiciones \\
Quiénes pueden ser citados --- art. 405 & partes, representantes, apoderados, personas jurídicas \\
Elección del absolvente --- art. 406 & empresa, conocimiento personal, oposición \\
Citación y domicilios & cédula, domicilio constituido vs. real, tres días \\
Reserva del pliego --- art. 410 & sobre cerrado, media hora antes, caducidad \\
Forma de las posiciones --- art. 411 & afirmativas, un hecho, personales, sí o no \\
Posición impertinente --- art. 414 & oposición sin fundamentación, efecto en sentencia \\
Confesión ficta --- art. 417 & incomparecencia, evasiva, consecuencia, plena prueba \\
Valor probatorio de la confesión ficta & doctrina dividida, fallo CSJN 2001 \\
Confesión expresa --- art. 423 & plena prueba, tres excepciones, comparación con ficta \\
Confesión extrajudicial --- art. 424 & policial, pericia, alcance, indivisibilidad \\
\bottomrule
\end{longtable}

% ============================================================
% CAPÍTULO 1 — PRUEBA TESTIMONIAL
% ============================================================
\chapter{La prueba testimonial}

\section{Concepto de testigo}

La prueba testimonial comienza a regularse a partir del artículo 426 del Código Procesal Civil y Comercial de la Nación (CPCCN). Se define al \textbf{testigo} como un tercero que no es parte en el proceso, al cual se convoca para que preste declaración sobre hechos de su conocimiento. El testigo es citado durante la etapa probatoria para que informe lo que sabe sobre los hechos controvertidos o ventilados en la causa.

\section{Objeto del testimonio y valor según el tipo de percepción}

El objeto del testimonio son los hechos alegados por las partes en sus escritos postulatorios. El testigo puede conocer esos hechos a través de distintas vías:

\begin{itemize}
    \item Haberlos visto directamente.
    \item Haberlos escuchado.
    \item Haberlos conocido a través del relato de otra persona.
\end{itemize}

No todos los testimonios poseen la misma eficacia probatoria: no es lo mismo quien presenció un hecho directamente que quien lo conoce de oídas, por lo que la \textbf{valoración} que realice el juez al dictar sentencia será distinta según el tipo de percepción. Sin perjuicio de ello, todos esos tipos de testimonio resultan admisibles.

\begin{formula}{Artículo 426 CPCCN --- Procedencia de la prueba testimonial}
Toda persona mayor de 14 años podrá ser propuesta como testigo y tendrá el deber de comparecer y declarar, salvo las excepciones establecidas por la ley.
\end{formula}

\section{Deberes del testigo}

Sobre todo testigo recaen tres deberes fundamentales: \textbf{comparecer}, \textbf{declarar} y \textbf{decir verdad}, cuyo alcance y excepciones se sintetizan en el siguiente cuadro.

\begin{table}[H]
\centering
\begin{tabularx}{\textwidth}{
    >{\raggedright\arraybackslash}p{0.16\textwidth}
    >{\raggedright\arraybackslash}X
    >{\raggedright\arraybackslash}X
}
\toprule
\textbf{Deber} & \textbf{Alcance} & \textbf{Excepciones} \\
\midrule
\textbf{Comparecer} & Es una carga pública. Toda persona convocada como testigo debe presentarse ante el juzgado. & Cargo que implique declarar por escrito; causa justificada debidamente acreditada; irregularidad en la notificación (menos de tres días de anticipación). \\
\textbf{Declarar} & El testigo tiene la obligación de prestar declaración testimonial. & Secreto profesional (militar, científico, artístico); declaración que perjudique el honor del testigo o lo exponga a enjuiciamiento penal. \\
\textbf{Decir verdad} & Antes de declarar, el testigo presta juramento o promesa de decir verdad. Se le informan las consecuencias penales del falso testimonio. & No hay excepciones al juramento, pero si el testigo miente incurre en el delito de falso testimonio. \\
\bottomrule
\end{tabularx}
\end{table}

\begin{important}
Respecto del deber de comparecer, resulta aplicable la denominada \textbf{regla de los 70 kilómetros}: si el domicilio del testigo se encuentra dentro de ese radio respecto del tribunal, tiene la obligación de concurrir personalmente. Si su domicilio está fuera de los 70 kilómetros, se lo convoca mediante la Ley 22.172, a través de un oficio dirigido al juzgado de turno en la sede de su domicilio. Ese oficio debe contener los datos personales del testigo y el interrogatorio completo, dado que la parte contraria también tiene derecho a controlar este medio de prueba y a formular sus propias preguntas.
\end{important}

\section{Testigos excluidos --- Artículo 427 y su vinculación con el Código Civil y Comercial}

El artículo 427 establece que no pueden ser ofrecidos como testigos los consanguíneos o afines en línea directa de las partes, ni el cónyuge aunque estuviere separado legalmente, salvo en los casos de reconocimiento de firmas. Los parientes colaterales (hermanos, tíos, sobrinos) sí pueden ser testigos.

\begin{formula}{Artículo 427 CPCCN --- Testigos excluidos}
No podrán ser ofrecidos como testigos los consanguíneos o afines en línea directa de las partes, ni el cónyuge aunque estuviere separado legalmente, salvo si se trata de reconocimiento de firmas.
\end{formula}

\begin{formula}{Artículo 711 del Código Civil y Comercial de la Nación --- Testigos en juicios de familia}
Los parientes y allegados a las partes pueden ser ofrecidos como testigos en juicios de familia. El fundamento es que, en los conflictos familiares, muchas veces quienes mejor conocen la problemática son precisamente los parientes o allegados, y excluirlos generaría una imposibilidad de probar los hechos alegados (por ejemplo, en casos de violencia familiar).
\end{formula}

\section{Los tres momentos de la prueba testimonial}

Todo medio probatorio se despliega en tres momentos:

\begin{enumerate}
    \item \textbf{Ofrecimiento}: se realiza en los escritos postulatorios (demanda y contestación de demanda).
    \item \textbf{Producción}: se lleva a cabo durante la etapa probatoria, una vez abierta la causa a prueba.
    \item \textbf{Valoración}: las partes valoran la prueba en los alegatos; el juez la valora al dictar sentencia.
\end{enumerate}

\section{Ofrecimiento de la prueba testimonial}

Para ofrecer un testigo en el escrito postulatorio, el código exige informar:

\begin{itemize}
    \item Nombre completo del testigo.
    \item Ocupación.
    \item Domicilio.
\end{itemize}

El artículo 333 del CPCCN agrega un elemento adicional: argumentar el objeto de la citación, es decir, explicar sintéticamente por qué se convoca a esa persona y sobre qué hechos declarará.

\begin{example}{Objeto de la citación}
Se ofrece a Juan Pérez porque estuvo presente en el momento del accidente, o porque es el médico que estuvo a cargo del tratamiento de la parte actora.
\end{example}

Pueden ofrecerse hasta ocho testigos. Si la complejidad del caso lo amerita, pueden admitirse más, previa decisión judicial sobre la necesidad de la ampliación.

\section{Producción de la prueba testimonial: citación a audiencia}

Una vez que la causa está abierta a prueba y el juez ha fijado las audiencias de testigos (en la etapa probatoria, luego de la audiencia preliminar), corresponde notificar a los testigos. Los medios de notificación admitidos son:

\begin{itemize}
    \item Cédula de notificación.
    \item Telegrama colacionado.
    \item Acta notarial.
    \item Carta documento.
\end{itemize}

La notificación debe cursarse con al menos tres días de anticipación a la audiencia. La cédula debe informar tanto la audiencia principal (primera) como la audiencia supletoria (segunda), con la advertencia de que si el testigo no concurre a la primera sin causa justificada, será convocado a la segunda con citación de la fuerza pública (es decir, por intermedio de la policía).

\section{Desarrollo de la audiencia testimonial}

Al comenzar la audiencia, el testigo se identifica con su documento nacional de identidad y brinda sus datos personales. Luego se le toma juramento o promesa de decir verdad, explicándole las consecuencias penales del falso testimonio. Inmediatamente después, se procede al interrogatorio preliminar.

\begin{formula}{Artículo 441 CPCCN --- Interrogatorio preliminar (generales de la ley)}
Aunque las partes no lo pidan, los testigos serán siempre preguntados por: 1) Su nombre, edad, estado, profesión y domicilio; 2) Si es pariente por consanguinidad o afinidad de alguna de las partes y en qué grado; 3) Si tiene interés directo o indirecto en el pleito; 4) Si es amigo íntimo o enemigo de alguna de las partes; 5) Si es dependiente, acreedor o deudor de algunos de los litigantes, o si tiene algún otro género de relación con ellos o algún interés en el pleito.
\end{formula}

Estas preguntas se denominan, en el lenguaje forense, ``generales de la ley''. Su propósito es evaluar si el testigo puede estar afectado por alguna de estas circunstancias, lo que podría incidir en la valoración de su testimonio e incluso en su idoneidad (conforme al art. 456 CPCCN).

\section{El interrogatorio: forma de formular las preguntas --- Artículo 443}

Luego del interrogatorio preliminar, cada parte puede formular su interrogatorio al testigo, el cual puede presentarse por escrito o realizarse a viva voz.

\begin{formula}{Artículo 443 CPCCN --- Forma de las preguntas}
Las preguntas no contendrán más de un hecho, serán claras y concretas. No se formularán las que estén concebidas en términos afirmativos, sugieran la respuesta o sean ofensivas o vejatorias. No podrán contener referencias de carácter técnico, salvo si fuesen dirigidas a personas especializadas.
\end{formula}

En la práctica forense, la pregunta se formula con la fórmula ritual: ``Para que jure, diga el testigo si sabe y cómo le consta\ldots [el hecho sobre el que se pregunta]''.

La pregunta debe ser abierta y no sugestiva, dado que el objetivo es conocer genuinamente lo que el testigo sabe y cómo lo sabe, y no confirmar una versión ya armada. Por ejemplo, si se pregunta si el testigo sabe qué sucedió tal día en tal lugar, se busca que el testigo relate libremente, y no que responda ``sí'' o ``no'' a una afirmación del abogado.

Si una pregunta viola los requisitos del artículo 443, la parte contraria puede oponerse. Esa oposición genera una incidencia dentro de la audiencia: si el juez hace lugar a la oposición, la pregunta queda rechazada y las costas de la incidencia se imponen a quien formuló la pregunta defectuosa.

\section{Testigo técnico vs. perito --- distinción conceptual}

El artículo 443, al mencionar las referencias técnicas, permite distinguir dos figuras que no deben confundirse.

\begin{table}[H]
\centering
\begin{tabularx}{\textwidth}{
    >{\raggedright\arraybackslash}p{0.24\textwidth}
    >{\raggedright\arraybackslash}X
    >{\raggedright\arraybackslash}X
}
\toprule
 & \textbf{Testigo técnico} & \textbf{Perito} \\
\midrule
\textbf{Fungibilidad} & No es fungible. Es una persona específica que conoció los hechos. & Es fungible. Se designa por sorteo de oficio por el juez. \\
\textbf{Conocimiento del hecho} & Conoció los hechos en el pasado, cuando ocurrieron. & Toma conocimiento del caso en el presente, al momento de su designación. \\
\textbf{Ejemplo} & El médico que trató al paciente (parte actora) durante su enfermedad. & El perito médico designado por sorteo para realizar la pericia médica en el proceso. \\
\bottomrule
\end{tabularx}
\end{table}

Al testigo técnico, por tratarse de una persona especializada, sí pueden formulársele preguntas de carácter técnico. Al perito, en cambio, se le formulan puntos de pericia en los escritos postulatorios.

\section{Negativa a responder --- Artículo 444}

\begin{formula}{Artículo 444 CPCCN --- Negativa a responder}
El testigo podrá rehusarse a contestar las preguntas si la respuesta lo expusiera a enjuiciamiento penal o comprometiera su honor, o si no pudiera responder sin revelar un secreto profesional, militar, científico, artístico o industrial.
\end{formula}

\section{Caducidad de la prueba testimonial --- Artículos 432 y 437}

\begin{important}
Corresponde distinguir la \textbf{negligencia} de la \textbf{caducidad} en la producción de la prueba. La negligencia se basa en un elemento subjetivo: el desinterés de la parte en producir su prueba a lo largo del período probatorio. La caducidad, en cambio, opera de forma objetiva: el código establece supuestos específicos que, al configurarse, producen automáticamente la pérdida del derecho a esa prueba.
\end{important}

\begin{formula}{Artículo 432 CPCCN --- Caducidad de la prueba testimonial}
A pedido de parte y sin sustanciación, se tendrá por desistido del testigo a la parte que lo propuso, en los siguientes supuestos: a) Si no hubiera activado la citación del testigo y este no hubiera comparecido por esta razón; b) Si no habiendo comparecido el testigo a la primera audiencia sin invocar causa justificada, no requiriera oportunamente las medidas de compulsión necesarias; c) Si fracasara la segunda audiencia por motivos no imputables a la parte y no solicitara nueva audiencia dentro del quinto día.
\end{formula}

\begin{formula}{Artículo 437 CPCCN --- Cuarto supuesto de caducidad}
Si la parte que ofreció al testigo no concurriera a la audiencia por sí o por apoderado y no hubiere dejado interrogatorio, se la tendrá por desistida de aquel sin sustanciación.
\end{formula}

\begin{table}[H]
\centering
\begin{tabularx}{\textwidth}{
    >{\raggedright\arraybackslash}p{0.06\textwidth}
    >{\raggedright\arraybackslash}p{0.28\textwidth}
    >{\raggedright\arraybackslash}X
    >{\raggedright\arraybackslash}p{0.16\textwidth}
}
\toprule
\textbf{N.º} & \textbf{Supuesto} & \textbf{Descripción} & \textbf{Norma} \\
\midrule
1 & No se notificó al testigo y no compareció & La parte no libró la cédula ni ningún medio de notificación al testigo. & Art. 432 inc. a) \\
2 & Testigo no compareció a la primera audiencia y no se solicitaron medidas de compulsión & El testigo faltó sin causa y la parte no gestionó la segunda audiencia con fuerza pública. & Art. 432 inc. b) \\
3 & Fracasó la segunda audiencia y no se pidió nueva dentro de los cinco días & La segunda audiencia fracasó por causas no imputables a la parte, pero esta no solicitó nueva audiencia dentro del quinto día. & Art. 432 inc. c) \\
4 & La parte oferente no concurrió y no dejó interrogatorio & El testigo sí compareció, pero la parte que lo ofreció no se presentó ni dejó el pliego de preguntas. & Art. 437 \\
\bottomrule
\end{tabularx}
\end{table}

\begin{note}
El desistimiento del testigo por quien lo ofreció es válido --sin necesidad de consentimiento de la contraparte-- siempre que sea previo a la celebración de la audiencia.
\end{note}

\section{Idoneidad del testigo --- Artículo 456}

\begin{formula}{Artículo 456 CPCCN --- Idoneidad de los testigos}
Dentro del plazo de prueba, las partes podrán alegar y probar acerca de la idoneidad de los testigos. El juez apreciará, según las reglas de la sana crítica y en oportunidad de dictar sentencia definitiva, las circunstancias y motivos que corroboren o disminuyan la fuerza de las declaraciones.
\end{formula}

Este artículo permite a las partes generar un incidente por separado del proceso principal para demostrar que un testigo no era idóneo para declarar, pudiendo ofrecerse en ese incidente toda la prueba que se considere necesaria. Se trata de un mecanismo de defensa frente a una declaración que, si bien era formalmente admisible, reveló en su contenido elementos que cuestionan la imparcialidad o credibilidad del testigo.

% ============================================================
% CAPÍTULO 2 — PRUEBA DE CONFESIÓN
% ============================================================
\chapter{La prueba de confesión}

\section{Concepto de confesión y distinción entre confesión espontánea y provocada}

\begin{definition}{Confesión}
La confesión, como medio de prueba, consiste en una declaración formulada por quien es parte en el proceso, sobre hechos personales o de su conocimiento, que produce un efecto desfavorable al confesante y favorable para la otra parte.
\end{definition}

Corresponde distinguir entre:

\begin{itemize}
    \item \textbf{Confesión espontánea}: ocurre cuando la parte demandada reconoce un hecho en la contestación de demanda. Al reconocer ese hecho, genera plena prueba sobre él, sin necesidad de ningún otro medio probatorio.
    \item \textbf{Confesión provocada (absolución de posiciones)}: es el medio de prueba propiamente dicho. Se produce en audiencia, donde la parte absuelve las posiciones formuladas por la contraria.
\end{itemize}

La confesión en sí no constituye el medio de prueba, sino que es la absolución de posiciones --la confesión provocada en juicio-- el medio de prueba propiamente dicho.

\section{Momento del ofrecimiento y producción}

Al igual que la prueba testimonial, la confesión se ofrece en los escritos postulatorios. La producción de este medio de prueba --la absolución de posiciones-- se lleva a cabo en la audiencia preliminar, conforme al artículo 360 inciso 4 del CPCCN. En la práctica, muchos jueces fijan una audiencia posterior específicamente para este fin, aunque la norma establece la audiencia preliminar como el momento de producción.

\section{Quiénes pueden ser citados a absolver posiciones --- Artículo 405}

\begin{formula}{Artículo 405 CPCCN --- Personas que pueden ser citadas}
Podrán asimismo ser citados a absolver posiciones, además del actor y el demandado: los representantes de los incapaces por los hechos en que hayan intervenido personalmente en ese carácter; los apoderados por hechos realizados en nombre de sus mandantes, estando vigente el mandato, y por hechos anteriores cuando estuvieran sus representados fuera del lugar en que se sigue el juicio, siempre que el apoderado tuviese facultades para ello y la parte contraria lo consienta; y los representantes legales de las personas jurídicas, sociedades o entidades colectivas que tuviesen facultad para obligarlas.
\end{formula}

\section{Elección del absolvente en personas jurídicas --- Artículo 406}

Para las empresas y personas jurídicas, el artículo 406 permite oponerse, dentro del quinto día de notificada la audiencia, a que absuelva el representante elegido por el ponente, siempre que:

\begin{enumerate}
    \item Se alegue que ese representante no intervino personalmente ni tuvo conocimiento directo de los hechos.
    \item Se indique en el mismo escrito el nombre del representante que absolverá las posiciones.
    \item Se deje constancia de que esa persona fue notificada de la audiencia y presta su conformidad.
\end{enumerate}

\section{Citación: domicilio constituido vs. domicilio real}

La notificación de la audiencia de absolución de posiciones debe hacerse por cédula. Corresponde distinguir según la forma de actuación de la parte en el proceso.

\begin{table}[H]
\centering
\begin{tabularx}{\textwidth}{
    >{\raggedright\arraybackslash}p{0.22\textwidth}
    >{\raggedright\arraybackslash}X
    >{\raggedright\arraybackslash}p{0.24\textwidth}
}
\toprule
\textbf{Modalidad} & \textbf{Tipo de actuación} & \textbf{Domicilio donde se notifica} \\
\midrule
Parte con letrado patrocinante (actúa por derecho propio) & Firma sus presentaciones junto con su abogado patrocinante. & Domicilio constituido (generalmente el estudio del abogado). \\
Parte representada por apoderado & El apoderado actúa en nombre de la parte. & Domicilio real de la parte (no el domicilio procesal del apoderado). \\
\bottomrule
\end{tabularx}
\end{table}

\begin{important}
La cédula debe entregarse con al menos tres días de anticipación. En caso de urgencia debidamente justificada, el juez puede reducir ese plazo mediante resolución, pero la anticipación no puede ser inferior a un día.
\end{important}

\section{Reserva del pliego --- Artículo 410}

\begin{formula}{Artículo 410 CPCCN --- Reserva del pliego de posiciones}
El pliego deberá ser entregado en secretaría hasta media hora antes de la hora fijada para la audiencia, en sobre cerrado, al que se le pondrá cargo. Si la parte que pidió las posiciones no comparece sin justa causa a la audiencia, no hubiese dejado pliego, y comparece el citado, perderá el derecho a exigirlas.
\end{formula}

El sobre cerrado se presenta en la mesa de entradas del juzgado hasta media hora antes del inicio de la audiencia. El fundamento de este límite temporal es evitar que la parte ponente prepare dos pliegos distintos --uno para el caso de que el absolvente comparezca y otro para el caso de incomparecencia-- y elija el que más le convenga según las circunstancias.

\section{Forma de redactar las posiciones --- Artículo 411}

\begin{formula}{Artículo 411 CPCCN --- Forma de las posiciones}
Las posiciones serán claras y concretas, no contendrán más de un hecho, serán redactadas en forma afirmativa y deberán versar sobre puntos controvertidos que se refieran a la actuación personal del absolvente. El juez podrá modificar de oficio, sin recurso alguno, el orden y los términos de las posiciones propuestas por las partes sin alterar su sentido, y podrá asimismo eliminar las que fueran manifiestamente inútiles o superfluas.
\end{formula}

La diferencia fundamental con el interrogatorio de testigos radica en que, mientras las preguntas a los testigos deben ser abiertas y no sugestivas, las posiciones de la confesión deben formularse en términos afirmativos, de manera que el absolvente pueda responder simplemente ``sí'' o ``no'' (aunque luego puede agregar las aclaraciones que estime pertinentes). La fórmula ritual es: ``Para que jure como que es cierto que usted se encontraba en tal lugar tal día''. Eso constituye la posición.

\begin{important}
El ponente, al formular la posición, está reconociendo ese hecho. Por lo tanto: si el absolvente contesta ``sí, es cierto'', hay plena prueba sobre ese hecho (ambas partes lo reconocen); si el absolvente contesta ``no es cierto'', ese hecho deberá probarse por otro medio.
\end{important}

\section{Posición impertinente --- Artículo 414}

\begin{formula}{Artículo 414 CPCCN --- Posición impertinente}
La parte podrá negarse a contestar la posición si esta fuera impertinente. El juez valorará esa oposición al dictar sentencia y, si considera que la posición era pertinente, podrá tener por confeso al absolvente en ese punto.
\end{formula}

A diferencia de lo que ocurre en la prueba testimonial --donde la oposición a una pregunta genera una incidencia fundada--, en la confesión la oposición a una posición no requiere fundamentación: el abogado del absolvente simplemente manifiesta que se opone en los términos del artículo 414 y eso queda asentado en el acta. Sin embargo, si el juez, al dictar sentencia, considera que la posición era pertinente, tendrá por confeso al absolvente, lo que puede tener consecuencias muy significativas para el resultado del proceso.

\section{Interrogatorio de las partes --- Artículo 415}

\begin{formula}{Artículo 415 CPCCN --- Interrogatorio de oficio}
El juez podrá interrogar de oficio a las partes en cualquier estado del proceso, y estas podrán hacerse recíprocamente las preguntas y observaciones que juzgaren convenientes en la audiencia que corresponda, siempre que el juez no las declare superfluas o improcedentes.
\end{formula}

\begin{note}
Esta posibilidad --que las partes se formulen preguntas recíprocas en cualquier audiencia, o que el juez interrogue de oficio-- constituye un instrumento valioso para esclarecer la verdad de los hechos, aunque en la práctica se utiliza con poca frecuencia. Con la reforma procesal y el proceso por audiencias, este panorama debería modificarse.
\end{note}

\section{Confesión ficta --- Artículo 417}

\begin{formula}{Artículo 417 CPCCN --- Confesión ficta}
Si el citado no compareciera a declarar dentro de la media hora de fijada para la audiencia, o si habiendo comparecido rehusara a responder o respondiera de una manera evasiva, el juez al sentenciar lo tendrá por confeso sobre los hechos personales, teniendo en cuenta las circunstancias de la causa y las demás pruebas producidas. En caso de incomparecencia del absolvente, aunque no se hubiera extendido acta, se aplicará lo establecido en este párrafo, si el ponente hubiera presentado oportunamente el pliego de posiciones y el absolvente estuviere debidamente notificado.
\end{formula}

Para que se configure la confesión ficta deben reunirse los siguientes requisitos:

\begin{enumerate}
    \item El absolvente debe estar debidamente notificado de la audiencia (con todos los requisitos de forma ya explicados).
    \item El ponente debe haber presentado el pliego en sobre cerrado hasta media hora antes.
    \item El absolvente no comparece dentro de la media hora de la hora fijada, o se niega a contestar, o responde de manera evasiva.
\end{enumerate}

Configurados esos supuestos, el juez, al dictar sentencia, puede tener por confeso al absolvente. La norma aclara que el juez tomará en cuenta ``las circunstancias de la causa y las demás pruebas producidas'', es decir, no se trata de una consecuencia automática absoluta, sino que debe apreciarse en el contexto del proceso.

\section{Valor probatorio de la confesión ficta --- posiciones doctrinarias y jurisprudencia}

Existen distintas interpretaciones doctrinarias y jurisprudenciales sobre el valor de la confesión ficta.

\begin{table}[H]
\centering
\begin{tabularx}{\textwidth}{
    >{\raggedright\arraybackslash}p{0.24\textwidth}
    >{\raggedright\arraybackslash}X
}
\toprule
\textbf{Posición doctrinaria} & \textbf{Contenido} \\
\midrule
Primer sector & La confesión ficta importa plena prueba siempre que no haya otro elemento de juicio que la contraríe. \\
Segundo sector & La confesión ficta es plena prueba siempre que haya otros medios de prueba que colaboren y la corroboren. \\
Tercer sector & No tiene valor absoluto en ningún supuesto: siempre debe conjugarse la confesión ficta con el resto de las constancias y medios de prueba del expediente. \\
\bottomrule
\end{tabularx}
\end{table}

\begin{important}
\textbf{Fallo de la Corte Suprema de Justicia de la Nación --- 6 de febrero de 2001 (LL 2001-B, pág. 959).} La CSJN habilitó excepcionalmente el recurso extraordinario en una cuestión de prueba --materia normalmente ajena a la instancia extraordinaria por ser de derecho común-- y revocó la sentencia de cámara que había desestimado el valor de la confesión ficta. El caso involucraba a una empresa de transportes que había sido tenida por confesa en audiencia de absolución de posiciones, y existían además otros elementos de juicio que corroboraban esa conclusión.
\end{important}

No habría razón para atribuir mayor valor a una ficción legal --como la configuración de la confesión ficta-- que a la realidad que surge del resto de las constancias de la causa. Desde un enfoque sistémico del proceso, el juez, al dictar sentencia, debe valorar la totalidad de las constancias de la causa, sin que la confesión ficta resulte automáticamente superior a la realidad probatoria del expediente.

\section{Confesión expresa --- Artículo 423}

\begin{formula}{Artículo 423 CPCCN --- Efectos de la confesión judicial expresa}
La confesión judicial expresa constituirá plena prueba, salvo cuando: 1) dicho medio de prueba estuviere excluido por la ley respecto de los hechos que constituyen el objeto del juicio, o sobre derechos que el confesante no puede renunciar o transigir válidamente; 2) recayere sobre hechos cuya investigación prohíbe la ley; 3) se opusiere a las constancias de instrumentos fehacientes de fecha anterior agregados al expediente.
\end{formula}

La confesión judicial expresa constituye plena prueba, salvo en esos tres supuestos. Es decir, cuando existe confesión expresa no es necesario recurrir a la valoración del resto de la prueba (salvo las excepciones señaladas). Esta es la diferencia fundamental con la confesión ficta: la ficta es una creación legal que opera ante la inactividad del absolvente; la expresa es una declaración voluntaria y directa del confesante sobre un hecho.

\section{Confesión extrajudicial --- Artículo 424}

La confesión extrajudicial es aquella que se produce fuera del proceso judicial propiamente dicho. Puede tener lugar, entre otros supuestos:

\begin{itemize}
    \item En la declaración ante autoridad policial, que tiene valor y puede hacerse valer en el proceso.
    \item En las manifestaciones ante el perito designado judicialmente: si la parte, durante la realización de una pericia (por ejemplo, una pericia médica), le hace saber al perito algún hecho relevante, esas manifestaciones también pueden tener valor confesional.
\end{itemize}

Esta última hipótesis fue objeto de debate, aunque se considera que no afecta el derecho de defensa, porque ocurre en el marco del proceso y todas las partes tienen posibilidad de controlarlo. En todo caso, el juez la valorará en la sentencia junto con el resto de los medios de prueba.

\begin{formula}{Artículo 424 CPCCN --- Alcance de la confesión}
En caso de duda, la confesión deberá interpretarse en favor de quien la hace. La confesión es indivisible, salvo cuando el confesante invoca hechos impeditivos, modificativos o extintivos, o absolutamente separables e independientes; o cuando las circunstancias calificativas expuestas por quien confiesa fueren contrarias a una presunción legal o inverosímiles, y las modalidades del caso hicieran procedente la divisibilidad.
\end{formula}

\section{Ausencia del absolvente por enfermedad y regla de los 300 kilómetros}

Si el absolvente no puede concurrir a la audiencia por enfermedad, el artículo 418 permite que el juez o un funcionario del juzgado investido de autoridad se traslade al domicilio o lugar donde se encuentre el absolvente para tomar la declaración. La enfermedad debe justificarse con la debida anticipación y por escrito (art. 419).

\begin{important}
En cuanto a la distancia, existe una regla análoga a la de los 70 kilómetros prevista para los testigos: si el absolvente domicilia dentro de los 300 kilómetros del juzgado, debe concurrir ante el juez de la causa. Si su domicilio está a más de 300 kilómetros, se aplica la Ley 22.172 y se libra un oficio al juzgado competente en la zona del domicilio del absolvente, adjuntando el pliego de posiciones.
\end{important}

\section{Desarrollo de la audiencia de absolución de posiciones}

Al momento de celebrarse la audiencia, se procede a la apertura del sobre que contiene el pliego de posiciones. Si ambas partes ofrecieron el medio de prueba, se absuelven primero las posiciones a una de ellas y luego a la otra. La audiencia puede ser grabada o puede labrarse un acta escrita. Las posiciones --que ya estaban en el sobre cerrado-- se leen al absolvente, quien las contesta por sí o por no, pudiendo agregar las aclaraciones que estime pertinentes.

% ============================================================
% CUADRO CORNELL DE REPASO
% ============================================================
\chapter*{Cuadro Cornell de repaso --- Prueba testimonial y prueba de confesión}
\addcontentsline{toc}{chapter}{Cuadro Cornell de repaso}

\begin{longtable}{
    >{\raggedright\arraybackslash}p{0.27\textwidth}
    >{\raggedright\arraybackslash}p{0.58\textwidth}
}
\toprule
\textbf{Preguntas / palabras clave} & \textbf{Notas / respuestas} \\
\midrule
\endfirsthead
\toprule
\textbf{Preguntas / palabras clave} & \textbf{Notas / respuestas} \\
\midrule
\endhead

\textbf{¿Quién es un testigo?} &
Es un tercero ajeno al proceso, convocado en la etapa probatoria para declarar sobre hechos de su conocimiento a través de sus sentidos (vista, oído) o del relato de terceros. No es parte del proceso. \\

\textbf{¿Cuáles son los deberes del testigo?} &
Tres deberes: 1) comparecer (carga pública; se puede notificar con fuerza pública en caso de incumplimiento); 2) declarar (salvo secreto profesional o riesgo de autoincriminación); 3) decir verdad (juramento previo; el falso testimonio es delito penal). \\

\textbf{¿Quiénes están excluidos como testigos?} &
Conforme al art. 427, los consanguíneos o afines en línea directa y el cónyuge. Excepción: en juicios de familia (art. 711 CCyC) se admiten parientes y allegados porque son quienes mejor conocen la problemática familiar. \\

\textbf{¿Cómo se ofrece la prueba testimonial?} &
En los escritos postulatorios, indicando nombre completo, ocupación y domicilio del testigo, más el objeto de la citación (art. 333). Se pueden ofrecer hasta ocho testigos; el juez puede admitir más si es necesario. \\

\textbf{¿Qué son las ``generales de la ley''?} &
Son las preguntas del interrogatorio preliminar (art. 441), que se formulan a todo testigo antes del interrogatorio de fondo: datos personales, parentesco con las partes, interés en el pleito, amistad o enemistad, relación de dependencia. Permiten evaluar la idoneidad e imparcialidad del testigo. \\

\textbf{¿Cómo deben formularse las preguntas a un testigo?} &
Deben: contener un solo hecho; ser claras y concretas; no estar formuladas en términos afirmativos; no sugerir la respuesta; no ser vejatorias u ofensivas; no contener referencias técnicas (salvo al testigo técnico). Se formulan: ``Para que jure, diga el testigo si sabe y cómo le consta\ldots''. \\

\textbf{¿Qué diferencia hay entre testigo técnico y perito?} &
El testigo técnico conoció los hechos en el pasado (ej.: médico tratante). El perito es fungible, designado por sorteo de oficio, toma conocimiento en el presente y dictamina sobre puntos de pericia formulados en los escritos postulatorios. Al primero se le pueden hacer preguntas técnicas; al segundo, no en esta etapa. \\

\textbf{¿Cuáles son los supuestos de caducidad de la prueba testimonial?} &
1) No se notificó al testigo (art. 432 a); 2) no se gestionaron medidas de compulsión tras la primera audiencia fallida (art. 432 b); 3) no se pidió nueva audiencia en cinco días tras fracasar la segunda (art. 432 c); 4) la parte oferente no concurrió ni dejó interrogatorio (art. 437). Los tres primeros son objetivos; todos operan a pedido de parte. \\

\textbf{¿En qué consiste la confesión como prueba?} &
En una declaración de quien es parte sobre hechos personales o de su conocimiento, con efecto desfavorable al confesante y favorable a la contraria. Puede ser espontánea (en la contestación de demanda) o provocada (absolución de posiciones en audiencia). \\

\textbf{¿Cómo se redactan las posiciones?} &
En forma afirmativa (``Para que jure como que es cierto que usted se encontraba en\ldots''), con un solo hecho, sobre actuación personal del absolvente, de manera que pueda contestarse por sí o por no. El ponente, al formularlas, reconoce el hecho. \\

\textbf{¿Qué es la confesión ficta?} &
Es la ficción legal del art. 417: si el absolvente, debidamente notificado, no comparece, se niega a contestar o responde evasivamente, el juez, al dictar sentencia, puede tenerlo por confeso en los hechos de las posiciones. No es automática: el juez evalúa las demás constancias de la causa. \\

\textbf{¿Qué valor probatorio tiene la confesión ficta frente a la confesión expresa?} &
La confesión judicial expresa constituye plena prueba (art. 423), salvo tres excepciones legales. La confesión ficta tiene valor discutido: para un sector doctrinal es plena prueba si no hay elementos en contra; para otro, requiere corroboración de otros medios; para un tercero, siempre debe cotejarse con el resto de la prueba. La CSJN (fallo del 6/2/2001, LL 2001-B, p. 959) revocó una sentencia que había desestimado la confesión ficta cuando existían otros elementos corroborantes. \\

\textbf{¿Cuándo opera la caducidad en la prueba de confesión?} &
Cuando el ponente no comparece a la audiencia, no deja el pliego, y el absolvente sí comparece: el ponente pierde el derecho a exigir la absolución de posiciones (art. 410 in fine). \\

\midrule

\multicolumn{2}{p{\textwidth}}{
\textbf{Resumen:}
Esta unidad analizó dos medios de prueba esenciales del proceso civil y comercial argentino, siguiendo en ambos casos una estructura trifásica común: ofrecimiento (en los escritos postulatorios), producción (en la etapa probatoria) y valoración (en alegatos y sentencia). La prueba testimonial convoca a terceros ajenos al proceso para que declaren sobre hechos de su conocimiento, imponiéndoles los deberes de comparecer, declarar y decir verdad bajo juramento. El régimen del CPCCN regula minuciosamente el modo de formular las preguntas (art. 443), el interrogatorio preliminar (art. 441) y los supuestos de caducidad (arts. 432 y 437). La prueba de confesión, en cambio, provoca una declaración de las propias partes sobre hechos personales: puede ser espontánea (reconocimiento en la contestación de demanda, que genera plena prueba de inmediato) o provocada a través de la absolución de posiciones, cuya forma afirmativa y específica distingue a este medio de la testimonial. La confesión ficta (art. 417) opera ante la incomparecencia o evasiva del absolvente notificado, y si bien genera una presunción legal desfavorable, su valoración debe conjugarse con el resto de las constancias de la causa, tal como lo señaló la Corte Suprema en el fallo del 6 de febrero de 2001. En conjunto, ambos medios de prueba ilustran la tensión permanente del proceso civil entre el formalismo procedimental y la búsqueda de la verdad objetiva de los hechos.
} \\

\bottomrule
\end{longtable}

\end{document}
